% !TeX program = xelatex
\documentclass[12pt,a4paper]{extarticle}

% 1. Подключение общей преамбулы (пакеты)
\input{../common/preamble}

% 2. Определение переменных для конкретного документа
\newcommand{\docname}{Source Listing}
\newcommand{\doccode}{\hseDocCodeBase\ 12 01-1}
\newcommand{\doccodelu}{\hseDocCodeBase\ 12 01-1-ЛУ}

% 3. Подключение общих команд (проект, студент, руководитель)
\input{../common/commands}

% 4. Подключение стилей (колонтитулы, заголовки)
\input{../common/styles}

\begin{document}
\sloppy

% 5. Генерация титульных листов по шаблону
\input{../common/title_template}


% ============================================================
%  ENGLISH SOURCE LISTING BODY (project.lang == en)
%  §2.3.4.2: an English applied project is documented in English.
%  ГОСТ 19.401-78 «Program Text» rendered in English; references
%  follow IEEE citation style (§C of the practice program).
% ============================================================

% ============================================
% ANNOTATION
% ============================================

\section*{ANNOTATION}
\addcontentsline{toc}{section}{ANNOTATION}

% Hint: list which information about the program is provided in the
% document (repository map, purpose of services, libraries, contracts, etc.).
This document «\projectname. Source Listing» provides information on the composition, hierarchical structure and software used by the system under development. The document contains a repository map, the purpose of the implemented components, a description of the shared libraries and a description of the infrastructure layer.

% Hint: the sentence about a private repository is optional — remove it
% if your repository is public.
The full source code of the program resides in a remote Git repository. \hseOptional{The repository is private; access to it is granted on request to review the composition and implementation of the program.}

\clearpage

% ============================================
% CONTENTS
% ============================================
\hypersetup{linkcolor=black}
\tableofcontents
\hypersetup{linkcolor=blue}
\thispagestyle{fancy}
\clearpage

% ============================================
% 1. INTRODUCTION
% ============================================
\section{INTRODUCTION}

\subsection{Program name}

The name of the program is «\projectname».

% The English and short names are inserted automatically from the project
% settings (fields «Work title in English» and «Short program name») —
% there is nothing to edit here.
The name of the program in English is «\hseProjectNameEn».

The short name of the program is «\hseProjectNameShort».

\subsection{Brief description of the program's application area}

% Hint: in 2-3 sentences describe the purpose of the program and its
% application area; state the architectural style (monolith, microservices,
% monorepository) and the key technologies.
The program «\projectname» serves the following purpose: \hseFill{brief description of the application area}. It implements mechanisms for \hseFill{list of the program's key functions}. The architectural style of the system is \hseFill{name the architectural style}, and the system and uses \hseFill{key technologies and protocols}.

\clearpage

% ============================================
% 2. PROGRAM TEXT
% ============================================
\section{PROGRAM TEXT}

% Hint: state that the full program text resides in the repository and, if
% necessary, give the address and access conditions.
% Hint: the sentence about a private repository is optional — remove it
% if your repository is public.
The full program text resides in a remote Git repository at the following address: \hseFill{Git repository address}. \hseOptional{The repository is private; access to it is granted on request to review the composition and implementation of the program.}

\subsection{Repository structure}

% Hint: describe the top-level directory structure of the project.
% List the purpose of the main directories of your repository.
\begin{hseExample}
The software system is organised on the principle of \hseFill{repository organisation principle}. The top-level directory structure is as follows:

\begin{itemize}[leftmargin=1.25cm]
    \item \texttt{src/} --- application source code: services, modules and background workers;
    \item \texttt{libs/} --- shared libraries, client SDKs and runtime components;
    \item \texttt{infra/} --- infrastructure management code, configurations and deployment utilities;
    \item \texttt{docs/} --- sources of the technical, product and operational documentation;
    \item \texttt{scripts/} --- auxiliary development, build and artifact-generation scripts;
    \item \texttt{tests/} --- integration and system test materials.
\end{itemize}
\end{hseExample}

\subsection{Description of the modules (services)}

% Hint: for each key module (service) state its purpose. Fill the table
% below with your own data.
\begin{hseExample}
The modules (services) of the program are split into domains: \hseFill{list of domains or groups of modules}. The purpose of the key modules is given in Table~\ref{tab:modules}.

% The table is typeset with \captionof (not \begin{table}[H]): floating
% environments inside the yellow example block are not allowed.
\begin{center}
\small
\captionof{table}{Purpose of the program modules}
\label{tab:modules}
\begin{tabularx}{\textwidth}{|>{\raggedright\arraybackslash}p{0.32\textwidth}|X|}
\hline
\textbf{Module} & \textbf{Purpose} \\
\hline
\texttt{\hseFill{module name}} & \hseFill{brief purpose of the module} \\
\hline
\texttt{\hseFill{module name}} & \hseFill{brief purpose of the module} \\
\hline
\end{tabularx}
\end{center}
\end{hseExample}

\subsection{Shared libraries}

% Hint: describe the shared (internal) libraries used and their purpose.
% If there are no shared libraries, this subsection may be removed.
\begin{hseExample}
To ensure development uniformity and code reuse, a set of internal libraries grouped by purpose is used.

\begin{center}
\small
\captionof{table}{Purpose of the shared-library groups}
\label{tab:libs}
\begin{tabularx}{\textwidth}{|>{\raggedright\arraybackslash}p{0.32\textwidth}|X|}
\hline
\textbf{Group} & \textbf{Composition and purpose} \\
\hline
\texttt{\hseFill{library group}} & \hseFill{composition and purpose of the group} \\
\hline
\texttt{\hseFill{library group}} & \hseFill{composition and purpose of the group} \\
\hline
\end{tabularx}
\end{center}
\end{hseExample}

\subsection{API contracts and client interaction}

% Hint: describe the external and inter-service interaction (REST, gRPC,
% data formats, location of contract schemas).
\begin{hseExample}
External interaction with the program is described through \hseFill{external API protocols}. The endpoints accept and return data in the \hseFill{data format} format used by client systems. Inter-service interaction is built on \hseFill{inter-module interaction mechanism}: the source contract-schema files are stored in the \texttt{\hseFill{schema directory}} directory, and the client and server stubs are generated by \hseFill{code-generation tool}.
\end{hseExample}

\subsection{Data storage and event-driven interaction}

% Hint: state the DBMS, caching and messaging tools used, as well as the
% storage and delivery patterns applied.
\begin{hseExample}
The server-side components are implemented mainly in \hseFill{specify the languages and frameworks}. For storing business data \hseFill{DBMS used} is used, for caching and short-lived state --- \hseFill{caching tool}, and for asynchronous exchange of events --- \hseFill{message broker}.

% Hint: describe how reliable delivery of events and data consistency are
% ensured (e.g. Transactional Outbox/Inbox); if there is no event exchange,
% this paragraph may be removed.
Reliable delivery of events is ensured by \hseFill{event-delivery patterns applied}: the domain change and the event record are performed within a single transaction, after which a background process publishes the event to the \hseFill{message broker}.
\end{hseExample}

\subsection{Infrastructure and deployment}

% Hint: briefly describe the infrastructure layer: containerisation,
% orchestration tools, build and deployment automation.
\begin{hseExample}
The infrastructure layer automates deployment and operation:

\begin{itemize}[leftmargin=1.25cm]
    \item \textbf{\hseFill{orchestration tools directory}} --- manifests, configurations and operational scripts for \hseFill{infrastructure components};
    \item \textbf{scripts/} --- CI/CD and operational scripts;
    \item \textbf{docker-compose.yml} --- local infrastructure environment for development and debugging.
\end{itemize}
\end{hseExample}

\subsection{Contracts and documentation}

% Hint: describe where the contract schemas and project documentation are
% stored; if the subsection is not relevant, it may be removed.
\begin{hseExample}
\begin{itemize}[leftmargin=1.25cm]
    \item \textbf{\hseFill{contract-schema directory}} --- a hierarchical file structure defining the interfaces of \hseFill{list of the main interfaces};
    \item \textbf{docs/} --- a structured knowledge base split by domains: \hseFill{list of documentation domains};
    \item \textbf{README.md} and per-service \textbf{README.md} --- instructions on running, developing, the API composition and the operational specifics of individual components.
\end{itemize}
\end{hseExample}

\clearpage
\pagestyle{appendix}

% ============================================
% 3. REFERENCES
% ============================================
\section*{REFERENCES}
\addcontentsline{toc}{section}{REFERENCES}

% Hint: provide the list of references. Below are the mandatory ESPD
% normative documents and one software-source example — replace or
% extend them with your own. References follow IEEE citation style.
\renewcommand{\refname}{}
\begin{thebibliography}{99}
\vspace{-1cm}

\bibitem{python_docs}
Python Software Foundation, \enquote{Python documentation,} Accessed: Jan.~1, 2026. [Online]. Available: \url{https://docs.python.org/3/}

\bibitem{gost19101}
\emph{GOST 19.101-77. Unified System of Program Documentation. Types of Programs and Program Documents}. Moscow, Russia: Standards Publishing House, 2001.

\bibitem{gost19105}
\emph{GOST 19.105-78. Unified System of Program Documentation. General Requirements for Program Documents}. Moscow, Russia: Standards Publishing House, 2001.

\bibitem{gost19401}
\emph{GOST 19.401-78. Unified System of Program Documentation. Program Text. Requirements for Content and Presentation}. Moscow, Russia: Standards Publishing House, 2001.

\end{thebibliography}



\clearpage

% ============================================
% ЛИСТ РЕГИСТРАЦИИ ИЗМЕНЕНИЙ
% ============================================
\input{../common/change_log}

\end{document}
